\chapter{Methodology}
\label{ch:methodology}

This chapter describes the simulation and processing pipeline shared by the
experiments in \cref{ch:hyp1,ch:hyp2,ch:hyp3}. Each experiment varies the
scatterer/material configuration and the displacement or material change
under test, but reuses the same forward model, signal-conditioning steps,
migration implementations, and (from \cref{ch:hyp1} onward) the same
phase-plane estimator. The chapter closes with a short validation experiment
(\cref{sec:meth-resolution}) that establishes the amplitude-based resolution
floor of the migration algorithms before they are used to test any of the
three hypotheses.

\section{Forward Modelling with gprMax}

All B-scans are simulated with the open-source finite-difference
time-domain solver gprMax \citep{gprmax}. The source is a Ricker wavelet with
centre frequency $f_c=\SI{15}{\GHz}$ and time-zero offset
$t_0=\SI{0.0943}{\ns}$ (\cref{fig:res-setup}a), chosen so that its usable
bandwidth defines the dominant wavelength $\lambda$ used to express every
displacement scale in this thesis ($2\lambda$ down to $\tfrac{1}{32}\lambda$).
The computational domain is discretised on a uniform \SI{1}{mm} grid with
perfectly-matched-layer (PML) absorbing boundaries.
\draftnote{the figure titles encode the domain extent as ``4010~m''; the
same auto-titling code elsewhere strips decimal points from floats (e.g.\
a depth of \SI{0.676}{m} appears as \texttt{0676\_m}, and a position of
\SI{2.0}{m} appears as \texttt{20\_m}), so this almost certainly reads as a
domain of $\approx\SI{4.01}{m}$ rather than \SI{4010}{m} --- confirm against
the notebook before quoting a final value.} A zero-offset (collocated
transmitter and receiver) survey is simulated by sweeping a single
transmitter--receiver pair across the surface.

\section{Scatterer and Medium Models}

Two target geometries are used across the experiments:

\begin{itemize}
  \item \textbf{Point scatterers} (\cref{sec:meth-resolution,ch:hyp1}):
    perfect-electric-conductor (PEC) cylinders of radius $r=\SI{28}{mm}$,
    buried at a depth of \SI{0.676}{m} in ice. \Cref{sec:meth-resolution}
    places two such cylinders at a swept separation; \cref{ch:hyp1} instead
    holds one cylinder fixed as a baseline and displaces a second, in the
    lateral, vertical, or diagonal direction, by the same family of
    sub-wavelength steps.
  \item \textbf{A distributed fluid front} (\cref{ch:hyp2}): a
    sub-wavelength fracture whose air-filled and water-filled regions are
    separated by a front that advances laterally between baseline and
    monitor surveys, following the thin-layer reflectivity model of
    \cref{sec:th-material-change}.
\end{itemize}

\section{Signal Conditioning Pipeline}
\label{sec:meth-conditioning}

Every raw B-scan is processed identically before migration:
\begin{enumerate}
  \item \textbf{Background subtraction.} A background-only simulation (no
    scatterer) is subtracted trace-by-trace to suppress the direct
    air/ground wave and isolate the scatterer reflection
    (e.g.\ \cref{fig:res-bscans}).
  \item \textbf{Tapering and $t_0$ alignment.} An exponential decay taper
    suppresses late-arriving energy, a cosine end-taper removes hyperbola
    tails at the edge of the migration aperture, and a static shift aligns
    the surface reflection to $t=0$ (e.g.\ \cref{fig:res-taper}).
  \item \textbf{Noise injection (where stated).} \Cref{ch:hyp3} contaminates
    the conditioned B-scan with synthetic Laplace-distributed noise at
    $10\,\%$ of the signal standard deviation, fitted from real field data.
\end{enumerate}

\section{Migration Algorithms Implemented}

Every conditioned B-scan in \cref{sec:meth-resolution,ch:hyp1,ch:hyp2,ch:hyp3}
is migrated with all three algorithms derived in \cref{sec:th-migration}:
Kirchhoff delay-and-sum (PyLops zero-offset operator), Gazdag $f$-$k$
phase-shift migration, and gprMax-based time-reversal back-propagation. All
three share the implementation in \texttt{helper\_functions/migration.py}
(\texttt{PylopsKirchoffMigration}, \texttt{gazdag\_migration}, and
\texttt{write\_backprop\_files}) and the exploding-reflector convention
$\vmig=v/2$ of \cref{eq:vmig}, so that the same velocity model and the same
migration aperture are used for a baseline/monitor pair, which is required
for the displacement estimate of \cref{sec:meth-phaseplane} to be valid.

\section{The 2D Phase-Plane Shift-Estimation Pipeline}
\label{sec:meth-phaseplane}

\Cref{ch:hyp1,ch:hyp2} apply the theory of
\cref{sec:th-fourier-shift,sec:th-wls,sec:th-material-change} to the migrated
images produced above, via the function \texttt{estimate\_shift\_2d} defined
in \texttt{TimeLapse\_Processing.ipynb}. The end-to-end workflow is:
\begin{enumerate}
  \item \textbf{Migrate} the baseline and monitor B-scan with an identical
    velocity model and aperture (\cref{sec:th-migration}).
  \item \textbf{Crop a region of interest (ROI)} tightly around the target
    (in this thesis, a window of $\pm\,2.5\lambda$ about the scatterer or
    fracture apex, located on the Hilbert-envelope peak of the baseline
    image), so that static background structure elsewhere in the image
    cannot drag the fitted plane towards zero.
  \item \textbf{Taper, then take the 2D FFT} of both cropped images and form
    the cross-spectrum $\XS$ of \cref{eq:cross-spectrum-def}.
  \item \textbf{Mask and weight}: restrict the fit to $|\kz|,|\kx|<1.4\,k_{zc}$
    and weight every bin by $|\XS|$ (\cref{sec:th-mask-weight}).
  \item \textbf{Solve} the weighted least-squares system of
    \cref{eq:wls-preweighted} for $(\Dz,\Dx,c)$.
\end{enumerate}
For the distributed fluid front of \cref{ch:hyp2}, step 2 instead slides
a small window along the known fracture geometry, and step 5 is read for its
intercept $c$ at each position to build a spatial profile of material
change (\cref{eq:intercept-material}), rather than a single displacement
estimate.

\paragraph{Practical considerations carried over from field application.}
Two points are worth flagging for any future extension of this pipeline to
real field surveys (\cref{ch:discussion}), even though the synthetic data in
this thesis sidesteps them by construction: the ROI window should be roughly
$3$--$4$ times the target pulse width (too small clips the signal under
tapering, too large admits background noise into the fit); and a genuine
change in soil/ice moisture between surveys changes the true wave velocity
$v$, which the estimator would otherwise misinterpret as a spurious vertical
shift $\Dz$ unless the velocity is recalibrated against a known static
reflector first.

\section{Validation: Amplitude Resolution Floor of the Migration Algorithms}
\label{sec:meth-resolution}

Before any of the three hypotheses can be tested with confidence, the
migration pipeline above is validated on a simple, well-understood baseline
problem: how closely can two \emph{stationary} point scatterers be spaced
before the migration algorithms above can no longer tell them apart? This
validation experiment uses two of the PEC cylinders described above,
illuminated by a zero-offset GPR B-scan, and migrated with all three
algorithms. Its result is the amplitude-based resolution floor against which
every displacement-detection result in \cref{ch:hyp1,ch:hyp2,ch:hyp3} is
later compared.

\Cref{fig:res-setup} shows the forward-model setup: the Ricker source
wavelet, the gprMax domain and grid, and the swept separation between two
PEC cylinder scatterers ($r=\SI{28}{mm}$, depth \SI{0.676}{m}), from
$2\lambda$ down to $\tfrac{1}{16}\lambda$.

\begin{figure}[htbp]
  \centering
  \begin{subfigure}[t]{0.32\textwidth}
    \includegraphics[width=\linewidth]{RES_001_Ricker_Wavelet_f_c__15_GHz_t0__0943_ns.png}
    \caption{Source wavelet}
  \end{subfigure}
  \hfill
  \begin{subfigure}[t]{0.32\textwidth}
    \includegraphics[width=\linewidth]{RES_002_Resolution_Study__Model_Geometry__domain_4010_m_Δx__1_mm_PML.png}
    \caption{Model geometry}
  \end{subfigure}
  \hfill
  \begin{subfigure}[t]{0.32\textwidth}
    \includegraphics[width=\linewidth]{RES_003_Scatterer_Positions__PEC_Cylinders__r__28_mm_depth__0676_m.png}
    \caption{Scatterer positions}
  \end{subfigure}
  \caption{Forward-model setup for the resolution validation: (a) the Ricker
    source wavelet ($f_c=\SI{15}{\GHz}$, $t_0=\SI{0.0943}{\ns}$); (b) the
    gprMax domain and grid; (c) the swept separation between the two PEC
    cylinder scatterers, $r=\SI{28}{mm}$, depth \SI{0.676}{m}.}
  \label{fig:res-setup}
\end{figure}

\Cref{fig:res-bscans} shows the simulated zero-offset B-scans before and
after background subtraction, and \cref{fig:res-taper} the effect of the
standard tapering and $t_0$-shift conditioning of \cref{sec:meth-conditioning}
on the $2\lambda$ separation scenario.

\begin{figure}[htbp]
  \centering
  \begin{subfigure}[t]{0.48\textwidth}
    \includegraphics[width=\linewidth]{RES_004_GPR_B-Scans__Background_and_Separation_Models.png}
    \caption{Background and separation models}
  \end{subfigure}
  \hfill
  \begin{subfigure}[t]{0.48\textwidth}
    \includegraphics[width=\linewidth]{RES_005_GPR_B-Scans__Background_Subtracted.png}
    \caption{Background-subtracted}
  \end{subfigure}
  \caption{Raw zero-offset B-scans for the resolution validation, before (a)
    and after (b) background subtraction.}
  \label{fig:res-bscans}
\end{figure}

\begin{figure}[htbp]
  \centering
  \begin{subfigure}[t]{0.48\textwidth}
    \includegraphics[width=\linewidth]{RES_006_Effect_of_Tapering_and_t0_Shift__2λ_dataset_single_trace.png}
    \caption{Single trace}
  \end{subfigure}
  \hfill
  \begin{subfigure}[t]{0.48\textwidth}
    \includegraphics[width=\linewidth]{RES_007_B-scan_effect_of_tapering_and_t0_shift__2λ_dataset.png}
    \caption{Full B-scan}
  \end{subfigure}
  \caption{Effect of tapering and the $t_0$ shift on the $2\lambda$ dataset:
    (a) a single representative trace; (b) the complete B-scan.}
  \label{fig:res-taper}
\end{figure}

\Cref{fig:res-kirchhoff,fig:res-gazdag,fig:res-backprop} show the migrated
image for every separation scenario, for Kirchhoff, Gazdag, and
back-propagation migration respectively, each with a zoomed view around the
true scatterer depth.

\begin{figure}[htbp]
  \centering
  \begin{subfigure}[t]{0.48\textwidth}
    \includegraphics[width=\linewidth]{RES_008_Kirchhoff_Migration__All_Datasets____f_c15_GHz____aperture40.png}
    \caption{All separation scenarios}
  \end{subfigure}
  \hfill
  \begin{subfigure}[t]{0.48\textwidth}
    \includegraphics[width=\linewidth]{RES_009_Kirchhoff_Migration_zoomed____f_c15_GHz____aperture40.png}
    \caption{Zoomed around scatterer depth}
  \end{subfigure}
  \caption{Kirchhoff migration of the resolution validation
    ($f_c=\SI{15}{\GHz}$, aperture $=40$ traces): (a) all separation
    scenarios; (b) zoomed view.}
  \label{fig:res-kirchhoff}
\end{figure}

\begin{figure}[htbp]
  \centering
  \begin{subfigure}[t]{0.48\textwidth}
    \includegraphics[width=\linewidth]{RES_010_Gazdag_Phase-Shift_Migration__All_Datasets____f_c15_GHz____z.png}
    \caption{All separation scenarios}
  \end{subfigure}
  \hfill
  \begin{subfigure}[t]{0.48\textwidth}
    \includegraphics[width=\linewidth]{RES_011_Gazdag_Phase-Shift_Migration_zoomed____f_c15_GHz.png}
    \caption{Zoomed around scatterer depth}
  \end{subfigure}
  \caption{Gazdag phase-shift migration of the resolution validation
    ($f_c=\SI{15}{\GHz}$): (a) all separation scenarios; (b) zoomed view.}
  \label{fig:res-gazdag}
\end{figure}

\begin{figure}[htbp]
  \centering
  \begin{subfigure}[t]{0.48\textwidth}
    \includegraphics[width=\linewidth]{RES_012_Back-Propagation_E__All_Datasets____focus_at_1906_ns.png}
    \caption{$\lVert\mathbf{E}\rVert$, all scenarios}
  \end{subfigure}
  \hfill
  \begin{subfigure}[t]{0.48\textwidth}
    \includegraphics[width=\linewidth]{RES_013_Back-Propagation_E_zoomed____focus_at_1906_ns.png}
    \caption{$\lVert\mathbf{E}\rVert$, zoomed}
  \end{subfigure}

  \vspace{0.6em}
  \begin{subfigure}[t]{0.48\textwidth}
    \includegraphics[width=\linewidth]{RES_014_Back-Propagation_Ez__All_Datasets____focus_at_1906_ns.png}
    \caption{$E_z$, all scenarios}
  \end{subfigure}
  \hfill
  \begin{subfigure}[t]{0.48\textwidth}
    \includegraphics[width=\linewidth]{RES_015_Back-Propagation_Ez_zoomed____focus_at_1906_ns.png}
    \caption{$E_z$, zoomed}
  \end{subfigure}
  \caption{Time-reversal back-propagation migration of the resolution
    validation, focused at $t=\SI{1906}{\ns}$: field-magnitude image (a,b)
    and the $E_z$ component (c,d), each with a zoomed view around the
    scatterer depth.}
  \label{fig:res-backprop}
\end{figure}

\Cref{fig:res-psf}a overlays the signed migrated amplitude from all three
algorithms at $f_c=\SI{15}{\GHz}$, and \cref{fig:res-psf}b plots the
normalised lateral point-spread function (PSF) extracted at the true
scatterer depth as a function of separation.

\begin{figure}[htbp]
  \centering
  \begin{subfigure}[t]{0.48\textwidth}
    \includegraphics[width=\linewidth]{RES_016_Migration_Comparison__Signed_Amplitude____f_c15_GHz____apert.png}
    \caption{Signed-amplitude comparison}
  \end{subfigure}
  \hfill
  \begin{subfigure}[t]{0.48\textwidth}
    \includegraphics[width=\linewidth]{RES_017_Normalised_Lateral_PSF_at_True_Scatterer_Depth.png}
    \caption{Normalised lateral PSF}
  \end{subfigure}
  \caption{(a) Signed migrated amplitude for Kirchhoff, Gazdag, and
    back-propagation migration overlaid at $f_c=\SI{15}{\GHz}$; (b) the
    normalised lateral point-spread function at the true scatterer depth,
    swept across separations from $2\lambda$ to $\tfrac{1}{16}\lambda$.}
  \label{fig:res-psf}
\end{figure}

All three migration algorithms collapse the two scatterer hyperbolae into
distinguishable amplitude peaks for separations down to roughly half a
wavelength, but the two peaks progressively merge into a single lobe as the
separation shrinks further: the amplitude image alone cannot certify two
scatterers, or a sub-wavelength displacement of one scatterer, below this
floor. \draftnote{state the precise separation at which the methods stop
resolving two distinguishable PSF peaks, read directly off
\cref{fig:res-psf}b, and comment on any difference between the three
algorithms.} This amplitude-based floor is the motivation for the
phase-plane approach developed in \cref{ch:theory} and tested against
Hypotheses 1--3 in \cref{ch:hyp1,ch:hyp2,ch:hyp3}.
